\chapter{CƠ SỞ LÝ THUYẾT VÀ NỀN TẢNG CÔNG NGHỆ}

\section{Tổng quan về công nghệ Blockchain}

\subsection{Định nghĩa và cấu trúc dữ liệu}

Blockchain, theo định nghĩa cơ bản nhất, là một cấu trúc dữ liệu dạng chuỗi liên kết các khối (block) theo trình tự thời gian, trong đó mỗi khối chứa một tập hợp các giao dịch đã được xác thực và một giá trị băm mật mã (cryptographic hash) tham chiếu đến khối liền trước. Cấu trúc này đảm bảo rằng bất kỳ thay đổi nào ở một khối trong quá khứ đều dẫn đến sự thay đổi giá trị băm lan truyền (cascade) đến tất cả các khối phía sau, khiến việc giả mạo dữ liệu trở nên không khả thi về mặt tính toán.

Về mặt hình thức, mỗi khối $B_i$ có thể được biểu diễn như một bộ bốn:

\begin{equation}
B_i = (H_{i-1}, \; T_i, \; nonce_i, \; timestamp_i)
\end{equation}

trong đó $H_{i-1} = \text{SHA-256}(B_{i-1})$ là giá trị băm của khối liền trước, $T_i = \{tx_1, tx_2, \ldots, tx_m\}$ là tập giao dịch trong khối, $nonce_i$ là giá trị được điều chỉnh bởi cơ chế đồng thuận, và $timestamp_i$ là nhãn thời gian tạo khối. Các giao dịch $T_i$ trong khối được tổ chức dưới dạng cây Merkle (Merkle Tree), cho phép xác minh sự tồn tại của một giao dịch cụ thể với độ phức tạp $O(\log m)$ mà không cần tải toàn bộ khối.

\subsection{Các đặc tính cốt lõi}

Bốn đặc tính cốt lõi tạo nên giá trị của công nghệ Blockchain bao gồm:

\begin{enumerate}
    \item \textbf{Tính bất biến (Immutability):} Dữ liệu một khi đã được ghi vào Blockchain gần như không thể bị sửa đổi hoặc xóa bỏ nhờ cơ chế liên kết hash giữa các khối và sự đồng thuận của mạng lưới phân tán. Muốn thay đổi một giao dịch trong khối $B_i$, kẻ tấn công phải tính lại hash của toàn bộ các khối $B_i, B_{i+1}, \ldots, B_n$ với tốc độ nhanh hơn cả phần còn lại của mạng --- một bài toán được chứng minh là không khả thi khi kẻ tấn công kiểm soát dưới 51\% năng lực tính toán.

    \item \textbf{Tính minh bạch (Transparency):} Trên các Blockchain công khai, mọi giao dịch đều có thể được kiểm tra và xác minh bởi bất kỳ ai, tạo nên một ``sổ cái'' mở (open ledger) cho toàn bộ hệ thống.

    \item \textbf{Tính phi tập trung (Decentralization):} Không có một thực thể đơn lẻ nào kiểm soát toàn bộ mạng lưới. Các bản sao dữ liệu được duy trì đồng thời trên hàng nghìn node trên toàn thế giới, loại bỏ điểm lỗi đơn (single point of failure).

    \item \textbf{Tính đồng thuận (Consensus):} Các node trong mạng phải đạt được sự thống nhất về trạng thái hợp lệ của chuỗi thông qua một giao thức đồng thuận, mà không cần tin tưởng lẫn nhau.
\end{enumerate}

\subsection{Cấu trúc dữ liệu: Chuỗi tuyến tính so với DAG}

Một khía cạnh quan trọng trong thiết kế Blockchain là lựa chọn cấu trúc dữ liệu nền tảng. Hai mô hình phổ biến nhất là chuỗi tuyến tính (Linear Chain) và đồ thị có hướng không chu trình (Directed Acyclic Graph --- DAG).

\textbf{Chuỗi tuyến tính (Linear Chain):} Là cấu trúc mặc định của Bitcoin và Ethereum. Các khối được sắp xếp theo trình tự tuyến tính $B_0 \rightarrow B_1 \rightarrow B_2 \rightarrow \ldots \rightarrow B_n$, mỗi khối chỉ có đúng một khối cha (parent block). Ưu điểm là đơn giản, dễ xác minh thứ tự giao dịch. Nhược điểm lớn nhất là tính tuần tự: tại mỗi thời điểm chỉ có một khối mới được chấp nhận, giới hạn nghiêm trọng thông lượng của hệ thống.

\textbf{Đồ thị DAG:} Avalanche X-Chain sử dụng cấu trúc DAG, trong đó mỗi giao dịch (hoặc đỉnh --- vertex) có thể tham chiếu đến nhiều giao dịch trước đó (nhiều cạnh cha), cho phép nhiều giao dịch không xung đột được xử lý song song. Về mặt lý thuyết đồ thị, DAG là một đồ thị có hướng $G = (V, E)$ thỏa mãn điều kiện không tồn tại chu trình --- tức không tồn tại đường đi $v_1 \rightarrow v_2 \rightarrow \ldots \rightarrow v_1$. Tính chất này đảm bảo rằng mối quan hệ ``xảy ra trước'' (happens-before) giữa các giao dịch luôn được bảo toàn.

\begin{table}[H]
\centering
\caption{So sánh cấu trúc dữ liệu Linear Chain và DAG}
\label{tab:chain_vs_dag}
\begin{tabular}{|l|p{5cm}|p{5.5cm}|}
\hline
\textbf{Tiêu chí} & \textbf{Linear Chain} & \textbf{DAG} \\ \hline
Cấu trúc & Chuỗi tuyến tính, mỗi block có 1 parent & Đồ thị có hướng, mỗi vertex có nhiều parent \\ \hline
Xử lý song song & Không (tuần tự từng block) & Có (các vertex không xung đột chạy đồng thời) \\ \hline
Thông lượng lý thuyết & Bị giới hạn bởi block time & Tùy thuộc vào mức song song hóa \\ \hline
Độ phức tạp xác minh & $O(n)$ duyệt chuỗi & $O(V + E)$ duyệt đồ thị BFS/DFS \\ \hline
Xác định thứ tự & Tự nhiên (vị trí trong chuỗi) & Cần sắp xếp topo (topological sort) \\ \hline
Sử dụng bởi & Bitcoin, Ethereum C-Chain & IOTA Tangle, Avalanche X-Chain \\ \hline
\end{tabular}
\end{table}

Avalanche tận dụng cả hai cấu trúc: X-Chain dùng DAG cho giao dịch tài sản (tốc độ cao), còn P-Chain và C-Chain dùng Linear Chain cho logic đồng thuận và Smart Contract (cần thứ tự xác định).

\subsection{Bài toán đồng thuận phân tán và Byzantine Fault Tolerance}

Bài toán đạt được sự đồng thuận trong một hệ thống phân tán có các thành viên có thể gặp sự cố hoặc hành xử không trung thực đã được hình thức hóa qua bài toán ``Các vị tướng Byzantine'' (Byzantine Generals Problem) bởi Lamport, Shostak và Pease vào năm 1982 \cite{lamport1982byzantine}.

Trong mô hình này, một nhóm $n$ tướng quân cần thống nhất kế hoạch tấn công hoặc rút lui, nhưng trong số đó có tối đa $f$ tướng phản bội có thể gửi thông điệp sai lệch. Kết quả quan trọng của bài báo chỉ ra rằng: hệ thống có thể đạt đồng thuận nếu và chỉ nếu $n \geq 3f + 1$, tức là số node trung thực phải chiếm hơn hai phần ba tổng số node.

Điều kiện này được gọi là ngưỡng chịu lỗi Byzantine (Byzantine Fault Tolerance --- BFT) và trở thành nền tảng lý thuyết cho mọi giao thức đồng thuận Blockchain từ đó đến nay.

\subsubsection{Giao thức PBFT (Practical Byzantine Fault Tolerance)}

Castro và Liskov \cite{castro1999practical} đề xuất giao thức PBFT năm 1999, là giao thức BFT thực tế đầu tiên. PBFT hoạt động qua ba pha tuần tự:

\begin{enumerate}
    \item \textbf{Pre-prepare:} Node leader (primary) gán số thứ tự cho yêu cầu client và phát multicast tới tất cả replica.
    \item \textbf{Prepare:} Mỗi replica xác nhận thứ tự đề xuất của leader, phát bản tin PREPARE. Khi một replica nhận được $2f$ bản tin PREPARE hợp lệ từ các replica khác, nó chuyển sang pha tiếp theo.
    \item \textbf{Commit:} Mỗi replica phát bản tin COMMIT. Khi nhận được $2f + 1$ bản tin COMMIT, replica thực thi yêu cầu và trả kết quả cho client.
\end{enumerate}

Tổng số bản tin trao đổi trong một vòng đồng thuận PBFT là $O(n^2)$, khiến giao thức này chỉ phù hợp cho mạng lưới nhỏ (dưới 100 node).

\subsubsection{Nakamoto Consensus (Proof-of-Work)}

Bitcoin sử dụng cơ chế đồng thuận dựa trên bằng chứng công việc (Proof-of-Work --- PoW). Mỗi node thi đua tìm giá trị $nonce$ sao cho:

\begin{equation}
H(B_{prev} \| T \| nonce) < \text{target}
\end{equation}

Node tìm được $nonce$ hợp lệ đầu tiên có quyền tạo khối mới và nhận phần thưởng. An toàn với giả định honest majority ($>$50\% hashrate trung thực), nhưng tiêu tốn năng lượng cực lớn --- theo Cambridge Centre for Alternative Finance \cite{cambridge2023bitcoin}, mạng Bitcoin tiêu thụ khoảng 120 TWh/năm (ước tính trung bình năm 2023, dải từ 67 đến 240 TWh), tương đương tiêu thụ điện của một số quốc gia.

\subsubsection{Proof-of-Stake (PoS) và Delegated PoS}

Proof-of-Stake thay thế năng lực tính toán bằng lượng tài sản stake. Validator được chọn tạo khối mới với xác suất tỷ lệ thuận với lượng token đã stake:

\begin{equation}
P(\text{validator}_i \text{ được chọn}) = \frac{w_i}{\sum_{j=1}^{n} w_j}
\end{equation}

với $w_i$ là trọng số stake của validator $i$. Cơ chế slashing (tịch thu stake) ngăn chặn hành vi bất lương: nếu validator ký hai khối khác nhau ở cùng một vị trí (equivocation), phần stake của họ bị tịch thu một phần hoặc toàn bộ.

Avalanche sử dụng biến thể PoS kết hợp Snowball: validator phải stake tối thiểu 2.000 AVAX (khoảng 50.000--100.000 USD tùy thời điểm) trên P-Chain để tham gia đồng thuận. Trọng số stake ảnh hưởng đến xác suất được chọn trong quá trình lấy mẫu Snowball.

\begin{table}[H]
\centering
\caption{So sánh các giao thức đồng thuận tiêu biểu}
\label{tab:consensus_comparison}
\begin{tabular}{|l|c|c|c|c|}
\hline
\textbf{Tiêu chí} & \textbf{PBFT} & \textbf{PoW} & \textbf{PoS} & \textbf{Snowball} \\ \hline
Ngưỡng BFT & $f < n/3$ & $f < n/2$ & $f < n/3$ stake & $f < n/2$ \\ \hline
Truyền thông & $O(n^2)$ & $O(n)$ & $O(n)$ & $O(k\log n)$ \\ \hline
Thông lượng & Thấp & Rất thấp & Trung bình & Cao \\ \hline
Finality & Giây & Phút--Giờ & Phút & $<$ 2 giây \\ \hline
Phi tập trung & Thấp & Cao & Trung bình & Cao \\ \hline
Năng lượng & Thấp & Rất cao & Thấp & Thấp \\ \hline
\end{tabular}
\end{table}


\section{Nền tảng Avalanche}

\subsection{Kiến trúc tổng quan}

Avalanche được thiết kế như một ``Nền tảng của các nền tảng'' (Platform of Platforms) với kiến trúc phân tầng linh hoạt. Thành phần cốt lõi là Primary Network --- một mạng lưới mặc định mà mọi validator phải tham gia --- bao gồm ba chuỗi tích hợp sẵn, mỗi chuỗi được tối ưu hóa cho một mục đích cụ thể \cite{avalabs2023docs}:

\begin{enumerate}
    \item \textbf{X-Chain (Exchange Chain):} Sử dụng cấu trúc dữ liệu DAG (Directed Acyclic Graph) thay vì chuỗi tuyến tính, chuyên xử lý các thao tác tạo và chuyển giao tài sản kỹ thuật số (digital assets). Giao thức đồng thuận Avalanche trên X-Chain cho phép nhiều đỉnh (vertex) trong DAG được xử lý song song, đạt thông lượng cực cao.

    \item \textbf{P-Chain (Platform Chain):} Là chuỗi quản trị siêu dữ liệu (metadata) của toàn bộ mạng lưới, chịu trách nhiệm:
    \begin{itemize}
        \item Quản lý danh sách validator và trọng số stake ($w_i$) của từng validator.
        \item Điều phối quá trình tạo và quản lý Subnet mới thông qua giao dịch \texttt{CreateSubnetTx}.
        \item Thêm validator vào Subnet qua giao dịch \texttt{AddSubnetValidatorTx}, hoặc ủy quyền stake qua \texttt{AddDelegatorTx}.
        \item Quản lý danh sách các Blockchain đang hoạt động trên mỗi Subnet.
    \end{itemize}

    \item \textbf{C-Chain (Contract Chain):} Triển khai máy ảo Ethereum (EVM), cho phép chạy Hợp đồng thông minh Solidity trên Avalanche. Đây là chuỗi được sử dụng rộng rãi nhất trong hệ sinh thái, tương thích hoàn toàn với các công cụ Ethereum (MetaMask, Hardhat, Remix, ethers.js).
\end{enumerate}

\begin{figure}[H]
\centering
\fbox{\parbox{0.85\textwidth}{\centering\vspace{2cm}\textit{[Hình 2.1: Kiến trúc tổng quan mạng lưới Avalanche --- Primary Network, X/P/C-Chain và các Subnet]}\vspace{2cm}}}
\caption{Kiến trúc tổng quan mạng lưới Avalanche}
\label{fig:avalanche_arch}
\end{figure}

\subsection{Giao thức đồng thuận Snowball}

Giao thức đồng thuận Snowball là thành phần cốt lõi tạo nên sự khác biệt của Avalanche so với các nền tảng Blockchain khác. Snowball thuộc họ giao thức ``Avalanche Consensus'' \cite{rocket2020avalanche}, được xây dựng theo nguyên lý leo thang từ đơn giản đến phức tạp qua bốn giao thức: Slush $\rightarrow$ Snowflake $\rightarrow$ Snowball $\rightarrow$ Avalanche.

\subsubsection{Giao thức Slush --- Nền tảng lấy mẫu}

Slush là giao thức đơn giản nhất, đặt nền móng cho toàn bộ họ giao thức. Ý tưởng cốt lõi dựa trên quan sát: nếu một nhóm người liên tục thăm dò ý kiến ngẫu nhiên lẫn nhau rồi cập nhật theo đa số, toàn bộ nhóm sẽ nhanh chóng hội tụ về một quyết định chung --- tương tự hiện tượng ``bầy đàn'' (herding) trong kinh tế học hành vi. Quy trình cụ thể:

\begin{enumerate}
    \item Mỗi node khởi tạo với một màu (quyết định): Đỏ hoặc Xanh.
    \item Tại mỗi vòng (round), node chọn ngẫu nhiên $k$ node từ mạng (bộ mẫu --- sample).
    \item Nếu hơn $\alpha \cdot k$ node trong bộ mẫu có cùng một màu (ngưỡng siêu đa số), node chuyển sang màu đó.
    \item Lặp lại trong $m$ vòng.
\end{enumerate}

Slush không có bộ nhớ (memoryless), nên có thể bị dao động. Các giao thức tiếp theo bổ sung cơ chế ``nhớ'' để khắc phục.

\subsubsection{Giao thức Snowflake --- Thêm bộ đếm tin cậy}

Snowflake cải tiến Slush bằng cách thêm một bộ đếm liên tiếp (consecutive counter) $cnt$. Mỗi khi kết quả truy vấn nhất quán với lựa chọn hiện tại, $cnt$ tăng thêm 1. Nếu không nhất quán, $cnt$ được reset về 0. Khi $cnt \geq \beta$ (ngưỡng quyết định), node coi quyết định là chung cuộc.

\subsubsection{Giao thức Snowball --- Thêm bộ đếm tin tưởng}

Snowball là phiên bản mạnh hơn Snowflake, bổ sung thêm bộ đếm tin tưởng (confidence counter) $d[c]$ cho mỗi lựa chọn $c$. Snowball tích lũy toàn bộ lịch sử phản hồi:

\begin{itemize}
    \item Mỗi truy vấn thành công cho màu $c$ $\Rightarrow$ $d[c] \leftarrow d[c] + 1$.
    \item Node luôn ưu tiên màu có $d[c]$ cao nhất, giúp tránh hiện tượng ``lật'' do nhiễu tạm thời.
    \item Quyết định chung cuộc khi bộ đếm liên tiếp $cnt \geq \beta$.
\end{itemize}

Thuật toán Snowball được mô tả bằng mã giả như sau:

\begin{lstlisting}[caption={Thuật toán Snowball (mã giả)}, label={lst:snowball}]
procedure Snowball(k, alpha, beta):
    color <- initial_preference
    d[Red] <- 0, d[Blue] <- 0  // confidence counters
    cnt <- 0                     // consecutive counter
    lastColor <- color

    while not decided:
        // Lay mau ngau nhien k node
        S <- random_sample(network, k)
        responses <- query(S, "What is your color?")

        // Dem so phieu cho moi mau
        for each c in {Red, Blue}:
            count[c] <- |{r in responses : r == c}|

        // Kiem tra sieu da so
        for each c in {Red, Blue}:
            if count[c] >= alpha * k:
                d[c] <- d[c] + 1
                if d[c] > d[color]:
                    color <- c
                if c == lastColor:
                    cnt <- cnt + 1
                else:
                    cnt <- 0
                    lastColor <- c

        if cnt >= beta:
            accept(color)
            decided <- true
\end{lstlisting}

\subsubsection{Giao thức Avalanche --- Mở rộng cho DAG}

Giao thức Avalanche (cùng tên với nền tảng) là phiên bản mạnh nhất, mở rộng Snowball cho cấu trúc DAG thay vì chỉ quyết định nhị phân. Trong Avalanche protocol, mỗi giao dịch $T$ được biểu diễn như một đỉnh (vertex) trong DAG, kèm theo tập tham chiếu ``parent'' tới các đỉnh trước đó. Cơ chế ``transitive voting'' cho phép một phiếu bầu cho đỉnh $v$ đồng thời bầu cho toàn bộ tổ tiên (ancestors) của $v$ trong DAG, giúp giảm đáng kể số lượng truy vấn cần thiết.

Giao thức duy trì một ``confidence'' cho toàn bộ DAG bằng cách: khi lấy mẫu $k$ node và hỏi ``Bạn có chấp nhận đỉnh $v$ không?'', câu trả lời ``có'' đồng nghĩa với việc chấp nhận tất cả đỉnh mà $v$ phụ thuộc. Nhờ đó, hàng trăm giao dịch có thể đạt đồng thuận chỉ trong vài vòng lấy mẫu.

\subsubsection{Phân tích toán học --- Mô hình chuỗi Markov}

Quá trình đồng thuận của Snowball có thể được mô hình hóa như một chuỗi Markov (Markov Chain) trên không gian trạng thái $(n_R, n_B)$, trong đó $n_R$ và $n_B$ lần lượt là số node đang chọn Đỏ và Xanh tại thời điểm hiện tại ($n_R + n_B = n$).

Gọi $P_{flip}(n_R, k, \alpha)$ là xác suất một node Đỏ chuyển sang Xanh sau một vòng lấy mẫu. Xác suất này được tính theo phân phối nhị thức:

\begin{equation}
P_{flip}(n_R, k, \alpha) = \sum_{j=\lceil \alpha k \rceil}^{k} \binom{k}{j} \left(\frac{n - n_R}{n}\right)^{j} \left(\frac{n_R}{n}\right)^{k-j}
\label{eq:pflip}
\end{equation}

Đây là xác suất mà khi lấy mẫu $k$ node, có ít nhất $\lceil \alpha k \rceil$ node thuộc phe Xanh. Khi $n_R$ đã chiếm đa số (ví dụ $n_R > 0.6n$), $P_{flip}$ giảm theo hàm mũ, tạo ra hiện tượng ``metastability''.

\textbf{Ví dụ minh họa bằng số:} Với $n = 1000$, $k = 20$, $\alpha = 0.8$ (yêu cầu $\geq 16/20$ phiếu):
\begin{itemize}
    \item Khi $n_R = 500$ (50-50): $P_{flip} \approx 0.006$ --- vẫn có khả năng dao động.
    \item Khi $n_R = 600$ (60-40): $P_{flip} \approx 2.3 \times 10^{-5}$ --- gần như không đổi.
    \item Khi $n_R = 700$ (70-30): $P_{flip} < 10^{-12}$ --- an toàn tuyệt đối.
\end{itemize}

Xác suất toàn hệ thống đảo ngược quyết định sau khi đạt $\beta$ vòng liên tiếp nhất quán có thể ước lượng bằng:

\begin{equation}
P_{reverse} \leq n \cdot (P_{flip})^{\beta}
\label{eq:preverse}
\end{equation}

Với $k=20$, $\alpha=0.8$, $\beta=20$, $n=2000$: $P_{reverse} < 2000 \times (10^{-12})^{20} = 2000 \times 10^{-240} \approx 0$, xác nhận tính an toàn gần tuyệt đối của giao thức.

\begin{table}[H]
\centering
\caption{Tham số Snowball tiêu biểu và ý nghĩa}
\label{tab:snowball_params}
\begin{tabular}{|c|c|p{7cm}|}
\hline
\textbf{Tham số} & \textbf{Giá trị mặc định} & \textbf{Ý nghĩa} \\ \hline
$k$ & 20 & Kích thước bộ mẫu mỗi vòng \\ \hline
$\alpha$ & 0.8 & Ngưỡng siêu đa số ($\geq$ 16/20 phiếu) \\ \hline
$\beta$ & 20 & Số vòng liên tiếp nhất quán để finalize \\ \hline
$n_{min}$ & 5 & Số validator tối thiểu trong Subnet \\ \hline
Staking min & 2.000 AVAX & Số token tối thiểu để trở thành validator \\ \hline
\end{tabular}
\end{table}

\subsection{Mô hình Subnet}

Subnet (Subnetwork) là khái niệm trung tâm của kiến trúc mở rộng Avalanche \cite{avalabs2023subnets}. Mỗi Subnet là một tập hợp con các validator trong mạng Avalanche, cùng thống nhất về trạng thái của một hoặc nhiều Blockchain riêng biệt.

\subsubsection{Cấu trúc và đặc điểm}

\begin{itemize}
    \item \textbf{Tập Validator riêng:} Mỗi Subnet tự quyết định bộ validator tham gia, có thể yêu cầu các điều kiện cụ thể (KYC, vị trí địa lý, phần cứng tối thiểu). Điều này rất quan trọng cho các ứng dụng cần tuân thủ quy định pháp lý (regulatory compliance).

    \item \textbf{Máy ảo tùy chỉnh:} Subnet có thể chạy bất kỳ VM nào --- Subnet-EVM (biến thể EVM), TimestampVM, hoặc VM tự phát triển bằng Go/Rust. Mỗi VM triển khai interface \texttt{block.ChainVM} với các phương thức: \texttt{BuildBlock()}, \texttt{ParseBlock()}, \texttt{SetPreference()}, \texttt{Accept()}, \texttt{Reject()}.

    \item \textbf{Tham số mạng độc lập:} Gas price, block size, block time, gas limit đều có thể tùy chỉnh riêng cho từng Subnet thông qua file Gensis JSON.

    \item \textbf{Chia sẻ bảo mật:} Các validator Subnet bắt buộc phải là validator của Primary Network (stake tối thiểu 2.000 AVAX), do đó gián tiếp được bảo vệ bởi lớp an ninh kinh tế của mạng chính.
\end{itemize}

\subsubsection{Quy trình tạo Subnet trên P-Chain}

Quy trình tạo một Subnet mới trên Avalanche P-Chain bao gồm các bước sau, minh họa bằng chuỗi giao dịch P-Chain:

\begin{enumerate}
    \item \textbf{CreateSubnetTx:} Người quản trị gửi giao dịch tạo Subnet mới, chỉ định tập ``control keys'' --- các khóa có quyền quản trị Subnet. P-Chain sinh SubnetID duy nhất.
    \item \textbf{CreateBlockchainTx:} Tạo một Blockchain mới chạy trên Subnet vừa tạo, chỉ định VM ID (ví dụ \texttt{srEXiWaHuhNyGwPUi444Tu47ZEDwxTWrbQiuD7FmgSAQ6X7Dy} cho Subnet-EVM) và Genesis Data.
    \item \textbf{AddSubnetValidatorTx} ($\times n$ lần): Thêm từng validator vào Subnet, chỉ định thời gian bắt đầu/kết thúc và trọng số stake.
    \item \textbf{Validator Sync:} Các validator node tải VM plugin, đồng bộ Genesis state, và bắt đầu tham gia đồng thuận trên Blockchain mới.
\end{enumerate}

Toàn bộ quy trình này khi thực hiện bằng CLI yêu cầu từ 30--60 phút cho người có kinh nghiệm. Hệ thống của em tự động hóa quy trình này xuống dưới 2 phút.

\begin{table}[H]
\centering
\caption{So sánh Avalanche Subnet với các mô hình App-Chain khác}
\label{tab:appchain_comparison}
\begin{tabular}{|l|p{3.2cm}|p{3.2cm}|p{3.5cm}|}
\hline
\textbf{Tiêu chí} & \textbf{Avalanche Subnet} & \textbf{Cosmos Zone} & \textbf{Polkadot Parachain} \\ \hline
Số lượng chuỗi & Không giới hạn & Không giới hạn & $\sim$100 slot \\ \hline
Chi phí tạo chuỗi & Thấp (phí P-Chain tx) & Trung bình (cần bộ Val riêng) & Rất cao (đấu giá slot) \\ \hline
Bảo mật chia sẻ & Có (qua Primary Network) & Không (tự vận hành) & Có (qua Relay Chain) \\ \hline
Tùy chỉnh VM & Hoàn toàn (bất kỳ VM) & Hoàn toàn (Cosmos SDK) & Giới hạn (WASM) \\ \hline
Interoperability & AWM & IBC Protocol & XCMP Protocol \\ \hline
Thời gian finality & $<$ 2 giây & 6--7 giây & 12--60 giây \\ \hline
\end{tabular}
\end{table}

\subsection{Máy ảo Ethereum (EVM) và Subnet-EVM}

Ethereum Virtual Machine (EVM) \cite{wood2014ethereum} là máy ảo ngăn xếp (stack-based virtual machine) thực thi bytecode của Hợp đồng thông minh. EVM là thành phần trung tâm tạo nên sức mạnh lập trình của Blockchain thế hệ 2 và 3. Kiến trúc EVM được mô tả chi tiết trong Ethereum Yellow Paper \cite{wood2014ethereum}, và sách ``Mastering Ethereum'' \cite{antonopoulos2018mastering} cung cấp hướng dẫn thực hành toàn diện.

\subsubsection{Kiến trúc EVM}

EVM hoạt động như một ``máy trạng thái toàn cầu'' (global state machine). Trạng thái thế giới (world state) $\sigma$ là một ánh xạ từ địa chỉ account (20 bytes) tới trạng thái account. Mỗi account bao gồm:
\begin{itemize}
    \item \textbf{Nonce:} Bộ đếm giao dịch (EOA) hoặc số contract đã tạo (Contract Account).
    \item \textbf{Balance:} Số dư token gốc (ETH trên Ethereum, AVAX trên C-Chain).
    \item \textbf{Storage Root:} Gốc cây Merkle Patricia Trie lưu trữ dữ liệu contract.
    \item \textbf{Code Hash:} Hash SHA3 của bytecode contract (rỗng nếu là EOA).
\end{itemize}

Khi một giao dịch gọi hàm contract, EVM thực thi bytecode theo từng opcode. Mỗi opcode tiêu tốn một lượng gas xác định. Nếu gas hết giữa chừng, toàn bộ thay đổi bị rollback (reverted). Cơ chế gas ngăn chặn vòng lặp vô hạn và đảm bảo mọi tính toán đều có chi phí.

\subsubsection{Subnet-EVM --- Biến thể cho Avalanche}

Subnet-EVM là biến thể EVM được Ava Labs phát triển riêng cho Avalanche Subnet, với nhiều tối ưu:
\begin{itemize}
    \item Cấu hình gas price, gas limit, block time linh hoạt qua Genesis JSON.
    \item Hỗ trợ precompiled contracts đặc thù: \texttt{FeeManager} (điều chỉnh phí động), \texttt{NativeMinter} (đúc token gốc), \texttt{AllowList} (whitelist/blacklist address).
    \item Tương thích 100\% với toolchain Ethereum: Solidity, ethers.js, web3.js, MetaMask, Hardhat.
    \item Tích hợp Avalanche Warp Messaging cho giao tiếp cross-Subnet.
\end{itemize}


\section{Công nghệ ảo hóa container --- Docker}

\subsection{Tổng quan kiến trúc Docker}

Docker là một nền tảng mã nguồn mở cho phép đóng gói ứng dụng cùng toàn bộ môi trường chạy (dependencies, libraries, system tools) vào một đơn vị triển khai nhỏ gọn gọi là Container \cite{docker2023docs}. Khác với máy ảo truyền thống (Virtual Machine --- VM) sử dụng hypervisor để giả lập toàn bộ phần cứng, Docker Container chia sẻ kernel của hệ điều hành chủ (host OS), do đó nhẹ hơn đáng kể về tài nguyên.

\begin{table}[H]
\centering
\caption{So sánh Docker Container và Virtual Machine}
\label{tab:docker_vs_vm}
\begin{tabular}{|l|p{5cm}|p{5.5cm}|}
\hline
\textbf{Tiêu chí} & \textbf{Docker Container} & \textbf{Virtual Machine} \\ \hline
Cách ly & Mức tiến trình (process-level) qua Linux Namespace & Mức phần cứng (hardware-level) qua Hypervisor \\ \hline
Kích thước & MB (không có OS riêng) & GB (bao gồm Guest OS) \\ \hline
Khởi động & Milli-giây đến giây & Phút \\ \hline
Hiệu năng & Gần native (chia sẻ kernel) & Overhead 5--15\% do ảo hóa \\ \hline
Quản lý tài nguyên & cgroups (giới hạn mềm/cứng) & Phân bổ cố định (vCPU, vRAM) \\ \hline
Mật độ & Hàng trăm container/host & Hàng chục VM/host \\ \hline
\end{tabular}
\end{table}

\subsection{Cơ chế cách ly của Linux Kernel}

Docker Container dựa trên hai cơ chế của Linux kernel:

\textbf{Linux Namespace:} Cung cấp cách ly trên 7 chiều:
\begin{itemize}
    \item \texttt{PID namespace}: Mỗi container có cây tiến trình riêng, PID 1 trong container là tiến trình chính (init process).
    \item \texttt{NET namespace}: Mỗi container có network stack riêng (interface, routing table, firewall rules).
    \item \texttt{MNT namespace}: Hệ thống file riêng biệt, ẩn file system của host.
    \item \texttt{UTS namespace}: Hostname và domain name riêng.
    \item \texttt{IPC namespace}: Hàng đợi message, shared memory segment cách ly.
    \item \texttt{USER namespace}: UID/GID mapping, cho phép root trong container không phải root trên host.
    \item \texttt{CGROUP namespace}: Cách ly view cgroups hierarchy.
\end{itemize}

\textbf{Control Groups (cgroups):} Giới hạn và theo dõi tài nguyên:
\begin{itemize}
    \item \texttt{cpu.shares}, \texttt{cpu.cfs\_quota\_us}: Giới hạn CPU theo tỷ trọng hoặc tuyệt đối.
    \item \texttt{memory.limit\_in\_bytes}: Giới hạn RAM hard limit; vượt quá sẽ bị OOM-kill.
    \item \texttt{blkio.weight}: Giới hạn I/O bandwidth cho disk.
    \item \texttt{pids.max}: Giới hạn số tiến trình tối đa trong container.
\end{itemize}

\subsection{Docker Networking --- Mô hình Bridge cho P2P}

Docker hỗ trợ nhiều driver mạng. Trong hệ thống của đề tài, em sử dụng Docker Bridge Network để kết nối các validator container:

\begin{enumerate}
    \item Docker daemon tạo một virtual bridge \texttt{docker0} (hoặc custom bridge) trên host.
    \item Mỗi container nhận một virtual ethernet interface (\texttt{veth pair}) kết nối vào bridge.
    \item Container nhận IP nội bộ từ subnet bridge (ví dụ \texttt{172.17.0.0/16}).
    \item Port mapping (\texttt{-p host\_port:container\_port}) cho phép truy cập từ bên ngoài.
\end{enumerate}

Đối với P2P giữa các validator, hệ thống cấp phát port tự động trong dải \texttt{19650--20650}, mỗi validator chiếm một cổng cho staking port (9651) và một cổng cho HTTP API (9650):

\begin{lstlisting}[caption={Cấu hình Docker Container cho Validator Node}, label={lst:docker_validator}, language={}]
docker run -d \
  --name avax-subnet-validator-1 \
  --network avalanche-net \
  -p 19650:9650 \     # HTTP API
  -p 19651:9651 \     # Staking port
  --cpus="2" \        # Gioi han 2 CPU cores
  --memory="4g" \     # Gioi han 4GB RAM
  -e AVAGO_NETWORK_ID=local \
  -e AVAGO_STAKING_PORT=9651 \
  avaplatform/avalanchego:latest
\end{lstlisting}

\subsection{Docker trong ngữ cảnh quản lý Validator Node}

Cách tiếp cận container hóa validator mang lại nhiều ưu điểm trong ngữ cảnh đề tài:

\begin{itemize}
    \item \textbf{Cách ly sự cố:} Nếu một validator crash hoặc bị tấn công, chỉ container đó bị ảnh hưởng. Các validator khác tiếp tục hoạt động bình thường nhờ PID và Network namespace riêng biệt.

    \item \textbf{Kiểm soát tài nguyên chặt chẽ:} Thông qua cgroups, hệ thống giới hạn CPU và RAM cho từng validator (ví dụ: 2 CPU, 4GB RAM mỗi node). Điều này ngăn ngừa tình trạng một validator tiêu tốn quá nhiều tài nguyên gây ảnh hưởng toàn hệ thống.

    \item \textbf{Triển khai nhanh:} Khởi động một validator container từ image AvalancheGo mất dưới 10 giây, so với cài đặt từ source code mất 5--10 phút.

    \item \textbf{Quản lý vòng đời qua API:} Docker Engine cung cấp REST API (v1.43) cho phép start, stop, remove, inspect container bằng code, không cần SSH vào server.
\end{itemize}


\section{Hệ thống giám sát phân tán}

\subsection{Prometheus --- Thu thập dữ liệu Time-Series}

Prometheus là hệ thống giám sát và cảnh báo mã nguồn mở, ban đầu được xây dựng tại SoundCloud và hiện là dự án thành phần thứ hai được tốt nghiệp (graduated) khỏi Cloud Native Computing Foundation (CNCF) sau Kubernetes \cite{prometheus2023docs}. Prometheus hoạt động theo mô hình Pull-based: chủ động truy vấn (scrape) dữ liệu từ HTTP endpoint \texttt{/metrics} của các mục tiêu theo chu kỳ cấu hình sẵn.

Dữ liệu được lưu trữ dưới dạng chuỗi thời gian (time-series), mỗi series được định danh bởi tên metric và tập nhãn (labels):

\begin{verbatim}
avalanche_c_handler_tx_processing_count{
  chain_id="2q9e4r6Mu3U68J7kR3MjNqfUyPMnx",
  instance="validator-1:9650"
}
\end{verbatim}

Prometheus hỗ trợ bốn loại metric cơ bản:

\begin{table}[H]
\centering
\caption{Bốn loại Prometheus Metric}
\label{tab:prometheus_metrics}
\begin{tabular}{|l|p{4cm}|p{5cm}|}
\hline
\textbf{Loại} & \textbf{Đặc điểm} & \textbf{Ví dụ trong hệ thống} \\ \hline
Counter & Chỉ tăng, không giảm & Tổng số giao dịch đã xử lý \\ \hline
Gauge & Giá trị tức thời, tăng/giảm & RAM usage hiện tại, số peer \\ \hline
Histogram & Phân phối giá trị theo bucket & Latency giao dịch (p50, p95, p99) \\ \hline
Summary & Phần trăm phân vị & Thời gian phản hồi API \\ \hline
\end{tabular}
\end{table}

\subsection{Grafana --- Trực quan hóa dữ liệu}

Grafana là nền tảng phân tích và trực quan hóa dữ liệu hàng đầu, kết nối với Prometheus qua ngôn ngữ truy vấn PromQL. Trong hệ thống, Grafana sẽ hiển thị dashboard với các panel: CPU/RAM usage per validator, block sync rate, transaction throughput, và peer connectivity.


\section{Công nghệ phát triển ứng dụng}

\subsection{NestJS --- Framework Backend}

NestJS là framework phát triển ứng dụng phía máy chủ cho Node.js \cite{nestjs2023docs}, lấy cảm hứng từ kiến trúc Angular với các khái niệm Module, Controller, Provider (Service) và Decorator. NestJS sử dụng TypeScript làm ngôn ngữ chính.

Kiến trúc NestJS tuân theo các nguyên tắc thiết kế quan trọng:

\begin{itemize}
    \item \textbf{Dependency Injection (DI):} Container IoC (Inversion of Control) quản lý vòng đời của tất cả Service. Khi một Controller cần \texttt{SubnetsService}, NestJS tự động tạo instance (singleton mặc định) và inject vào constructor. Cơ chế này giúp code dễ test (dùng mock) và dễ thay thế implementation.

    \item \textbf{Module System:} Mỗi chức năng được tổ chức thành module có Controllers, Providers, Imports, Exports riêng biệt. Module có thể import module khác để sử dụng Services exported. Trong hệ thống, \texttt{SubnetsModule} import \texttt{PrismaModule} để truy cập database.

    \item \textbf{Guards \& Interceptors:} Guards quyết định một request có được phép tiếp tục không (authentication, authorization). Interceptors can thiệp vào response (logging, caching, transformation). Pipes validate và transform input data.

    \item \textbf{WebSocket Gateway:} Hỗ trợ tích hợp Socket.IO \cite{socketio2023docs} để truyền dữ liệu giám sát thời gian thực qua kết nối WebSocket two-way (full-duplex).
\end{itemize}

\subsection{Next.js --- Framework Frontend}

Next.js là framework React phát triển bởi Vercel \cite{nextjs2023docs}, cung cấp:

\begin{itemize}
    \item \textbf{App Router:} Hệ thống routing dựa trên cấu trúc thư mục (file-system based routing) với hỗ trợ Server Components (render trên server, không gửi JavaScript) và Client Components (interactive trên trình duyệt).

    \item \textbf{Server-Side Rendering (SSR) và Static Generation (SSG):} Trang có thể được render sẵn (pre-rendered) từ server hoặc tại build time, cải thiện hiệu năng tải trang đầu (First Contentful Paint) và SEO.

    \item \textbf{React Server Components (RSC):} Cho phép render component trên server mà không gửi component code tới client, giảm đáng kể kích thước JavaScript bundle.

    \item \textbf{API Routes:} Cho phép tạo endpoint API nhẹ phía server, phù hợp cho proxy request hoặc xử lý logic trung gian.
\end{itemize}

\subsection{Prisma ORM --- Tương tác cơ sở dữ liệu}

Prisma là công cụ ORM thế hệ mới cho Node.js và TypeScript \cite{prisma2023docs}. Ba thành phần chính:

\begin{enumerate}
    \item \textbf{Prisma Schema:} File khai báo mô hình dữ liệu (\texttt{schema.prisma}) sử dụng cú pháp khai báo trực quan. Prisma tự động sinh migration SQL và Prisma Client từ schema.

    \item \textbf{Prisma Client:} API truy vấn type-safe --- nghĩa là khi viết \texttt{prisma.user.findUnique(\{where: \{email\}\})}, TypeScript kiểm tra tại compile-time rằng field \texttt{email} tồn tại trong model \texttt{User}. Điều này ngăn chặn lỗi runtime rất hiệu quả.

    \item \textbf{Prisma Migrate:} Công cụ quản lý phiên bản database schema. Trong development, lệnh \texttt{prisma db push} đồng bộ schema trực tiếp; trong production, \texttt{prisma migrate deploy} áp dụng migration files.
\end{enumerate}

\subsection{JSON Web Token (JWT) và OAuth 2.0}

\subsubsection{Cấu trúc JWT}

JWT (RFC 7519 \cite{rfc7519jwt}) là tiêu chuẩn truyền tải thông tin an toàn giữa các bên dưới dạng JSON. Cấu trúc gồm ba phần phân cách bởi dấu chấm:

\begin{verbatim}
xxxxx.yyyyy.zzzzz
 |       |      |
Header Payload Signature
\end{verbatim}

\begin{enumerate}
    \item \textbf{Header:} Chứa thuật toán ký (HS256) và loại token (\texttt{``typ'': ``JWT''}).
    \item \textbf{Payload:} Chứa các claim: \texttt{sub} (user ID), \texttt{email}, \texttt{iat} (issued at), \texttt{exp} (expiration).
    \item \textbf{Signature:} $\text{HMAC-SHA256}(\text{base64(header)} + \text{``.''} + \text{base64(payload)}, \; \text{secret})$
\end{enumerate}

Server xác minh token bằng cách tính lại signature và so sánh, không cần truy vấn database --- đây là ưu điểm lớn của JWT so với session-based authentication.

\subsubsection{Giao thức OAuth 2.0 --- Google Sign-In}

OAuth 2.0 (RFC 6749 \cite{rfc6749oauth}) là giao thức ủy quyền (authorization framework) cho phép ứng dụng bên thứ ba truy cập tài nguyên của người dùng trên một dịch vụ khác mà không cần biết mật khẩu. Trong hệ thống, em sử dụng Google OAuth 2.0 Authorization Code Flow:

\begin{enumerate}
    \item Client redirect tới Google Authorization Server kèm \texttt{client\_id} và \texttt{redirect\_uri}.
    \item Google hiển thị Consent Screen, người dùng đồng ý chia sẻ thông tin (email, tên).
    \item Google redirect về \texttt{callback\_url} kèm authorization code.
    \item Backend đổi code lấy access token từ Google Token Endpoint.
    \item Backend dùng access token gọi Google UserInfo API lấy email, name.
    \item Backend tìm hoặc tạo User trong database, phát hành JWT nội bộ.
\end{enumerate}

Ưu điểm của phương pháp này: người dùng không cần tạo mật khẩu mới, giảm rủi ro lộ credential, và tận dụng hệ thống bảo mật mạnh của Google (2FA, phát hiện đăng nhập bất thường).
